\documentclass[11pt]{article}
\usepackage[T1]{fontenc}
\usepackage{lmodern,microtype,amsmath,amssymb,amsthm,booktabs}
\usepackage[margin=1in]{geometry}
\usepackage[hidelinks]{hyperref}
\newtheorem{proposition}{Proposition}
\title{Complex rectangle sums improve the conditional saving beyond $2^{-31}$}
\author{Research extension prepared with OpenAI Codex}
\date{October 7, 2026}
\begin{document}
\maketitle
\begin{abstract}
Signed rectangle mixers reduce the complex network's side-role count to
$455434$ per invocation while preserving its absolute rank deficit.
Four explicit binary frame choices give nondegenerate, nonalternating
residuals in both circuit directions. Together with first/third-stage
scratch sharing and the retained bit primitive and compact-control assembly,
this yields the conditional witness
$\kappa=59/10^{11}=5.9\cdot10^{-10}>2^{-31}$ in
$T(n)=O(n(\log n)^{1-\kappa})$.
This is $4.72$ times the preceding $125/10^{12}$ witness and about $7.108$
times the original repository's $83/10^{12}$ witness.
The upstream theorem and previous compact-control extensions remain
assumptions. The new proof has not received independent review.
\end{abstract}

\section{Scope and progression}
The baseline is Douglas Colkitt's \texttt{CrocSwap/integer-mult-bounds},
commit \texttt{6e564879f51ae16f23d392e9e196c605f36d90df}, and its pinned
OpenAI manuscript, commit
\texttt{adc7f1241b42e322a6451854ab7e4b4c146bf78a}.
The model remains a fixed finite alphabet and a fixed finite number of
one-dimensional Turing-machine tapes. We retain the paired bit network at
$h_{\rm b}=50$, the compact-control movement and reservation proofs,
local exceptional repair, independent complex arity, generalized stopped
guard, Gaussian choice, resampling, and final rounding interfaces.
The latest extension changes the complex scalar circuit and its binary
frames, shares its scratch banks, and changes the final parameter choices.

Section 2 preserves the first scratch-sharing construction, which supported
$\kappa=125/10^{12}>2^{-33}$. Its disjoint-image matching is also used by
the stronger rectangle construction in Section 3. Section 4 supplies the
latest parameter witness. Finite tests and exact arithmetic support the
written proofs; they do not formally verify the multiplication theorem.

\section{Earlier stage-sharing construction}
\subsection{Sharing scratch in the complex network}
\label{sec:complex-stage-reuse}
This extension retains every scalar gate and every gate label of the original
complex network at $h=25$. It identifies certain first-stage scratch roles
with third-stage scratch roles. The bit network and all its interfaces remain
unchanged. The following proof supplies the binary residual obligation that
is additional to rational stage sharing.

Write $F=\mathbb F_2^h$, $\mathcal T=\binom{[h]}3$, $v=\binom h3$,
$m=h^3$, and $z=\binom{h-3}3+3(h-3)$. Work with the ordinary dot product
on $F$ and the induced product on $F^{\otimes3}$. Triple indicators have
norm one. Neighbors are distinct triples with even intersection.

\paragraph{A disjoint-image permutation.}
Identify $[h]$ with $\mathbb Z/h\mathbb Z$. Partition triples into translation
orbits. In each orbit choose its lexicographically first representative $T$
and the least positive shift $k$ with $T\cap(T+k)=\varnothing$; apply that
same shift to every member of the orbit. Such a shift exists for $h\ge8$,
because $T-T$ contains at most seven residues. Translation permutes the orbit,
even if it has a stabilizer. The resulting permutation $\pi$ of all triples
therefore satisfies $t_A\cdot t_{\pi(A)}=0$ for every $A$.

\paragraph{Two explicit role matchings.}
Fix a common ordering of triples, and write $n_S(j)$ for the $j$th neighbor
of $S$, $0\le j<z$. A first-stage side role with fixed second and third
triples $(A,B)$ and physical first-coordinate target/source pair
$(S,n_S(j))$ is identified with the third-stage side role having fixed first
and second triples $(S,\pi(A))$ and third-coordinate target/source pair
$(B,n_B(j))$. This is bijective: $A$ is recovered by $\pi^{-1}$, while
$S,B,j$ are read from the third-stage role. Separately, identify central role
$i\in[h]\cup\{*\}$ of the first-stage invocation $(A,B)$ with central
role $i$ of the third-stage invocation $(B,\pi(A))$. This is also bijective.
Second-stage scratch remains separate. No roles used simultaneously are
identified.

Each invocation restores arbitrary scalar scratch, so these identifications
preserve the signed exchange $(X,Y)\mapsto(-Y,X)$ and fix every remaining
scratch role. In particular the argument does not assume zero temporary values.

\paragraph{Nested binary frames.}
The last first-stage label on either type of shared role is
\[
 E=F\otimes\langle t_A\rangle\otimes\langle t_B\rangle,
 \qquad \dim E=h.
\]
For the paired side role the first third-stage label is
\[
 H_{\rm side}=
 \langle t_S\otimes t_{\pi(A)}\rangle^\perp_{F\otimes F}
 \otimes\langle t_B\rangle,
 \qquad \dim H_{\rm side}=h^2-1.
\]
For the paired central role it is
\[
 H_{\rm center}=
 \langle t_B\otimes t_{\pi(A)}\rangle^\perp_{F\otimes F}
 \otimes F,
 \qquad \dim H_{\rm center}=m-h.
\]
The disjointness of $A$ and $\pi(A)$ implies $E\subset H$ in both cases.
All these spaces are nondegenerate: their defining lines have norm one,
and they are tensor products of nondegenerate spaces. Hence the new edge
has nondegenerate orthogonal residual $R=H\cap E^\perp$.

Nondegeneracy alone would not suffice over $\mathbb F_2$. For a side bridge,
choose $x\notin S$ and $y\notin A$. The vector
$e_x\otimes e_y\otimes t_B$ belongs to $H_{\rm side}\cap E^\perp$
and has norm one. For a central bridge choose $x\notin B$, $y\notin A$
and any $q\in[h]$. Then $e_x\otimes e_y\otimes e_q$ belongs to
$H_{\rm center}\cap E^\perp$ and has norm one. Every new residual is
therefore nonalternating and admits an orthonormal basis by the original
binary-space lemma. The side and central residual dimensions are respectively
$h^2-1-h$ and $m-2h$.

All other edges retain their original residuals and orthonormal bases.
For nested labels $H=E\mathbin\perp R$, orthogonal weight additivity gives
\[
 q_H(x)-q_E(x)=\operatorname{wt}(P_Rx)
 =\sum_{u\in\mathcal B_R}\operatorname{wt}(u)[u\cdot x]\pmod4.
\]
Each basis vector has odd weight, so the new edge is implemented by exactly
$\dim R$ forward or inverse translation kernels, just as in the original
complex phase interface. This verifies the phase signs, not only the dimensions.

\paragraph{Exact deficit and endpoints.}
The old two edges on a paired role, $E\to F^{\otimes3}$ and $0\to H$,
have total residual dimension $m-h+\dim H$. Their replacement $E\to H$
has dimension $\dim H-h$, saving exactly $m$ per deleted role.
The total number of deleted roles is
\[
 R_0=v^2(vz+h+1)=19540339540000.
\]
All central decreases are retained, including those in distinct invocations
that now share a role. Thus, writing $N=v^3$ and $L=3v^2h(h+1)$,
\begin{align*}
 W_{\rm c}^{\rm share}&=2N+2v^2(vz+h+1)=39105013080000,\\
 s_{\rm c}^{\rm share}&=W_{\rm c}^{\rm share}m-2N+2L
                       =611015825672000000,\\
 W_{\rm c}^{\rm share}m-s_{\rm c}^{\rm share}&=3703000000,\\
 \eta_{\rm c}^{\rm share}&=\frac1{165005625}.
\end{align*}
Every surviving scratch role still starts with zero label, ends with the
full label, and is fixed by the scalar permutation. Data endpoints and their
weight-27 corrections are unchanged. The original signed-exchange correction
and tensor-column argument consequently give the same completed all-role
phase interface, with the displayed smaller $W,s$.

\begin{proposition}\label{prop:shared-complex-interface}
The shared complex network supports the original phase interface with
$m=15625$, $W=W_{\rm c}^{\rm share}$, $s=s_{\rm c}^{\rm share}$, and
$\sigma=1-627/10^{12}$.
\end{proposition}
\begin{proof}
The frame, residual, and endpoint arguments above prove the interface.
Exact logarithm enclosures certify
$\eta_{\rm c}^{\rm share}>(627/10^{12})\log(15625)$.
Since $e^{-x}>1-x$ for $x>0$, this gives $s/W<m^\sigma$.
\end{proof}

\paragraph{Depth and tape constants.}
No scalar gate is added, and the number of residual kernels decreases.
One can retain the original unshared $h=25$ node-depth bound
$E_{\rm old}=64(W_{\rm old}+m+1)^3$, since its gate schedule is unchanged
and bridges only delete boundary operations. The old $s_{\rm old}$ also
dominates $s_{\rm c}^{\rm share}$. Therefore the entire generalized guard
proof remains valid with $B_{\rm old}=s_{\rm old}+E_{\rm old}$ and its
original $C_0$. The layer recurrence uses the new $W,s$ for time and row
splitting, and these conservative old constants only for coefficient depth.
All matchings, bases, and schedules are fixed finite data, so the retained
fixed-number-of-tapes argument applies without an input-dependent resource.

\section{Stronger rectangle construction}
\subsection{Signed rectangle sums for the complex network}
\label{sec:complex-rectangles}
Retain $h=25$, $v=\binom h3=2300$, $m=h^3=15625$, and $N=v^3$.
The complex scalar side matrix has coefficient $+1/2$ on disjoint triples,
$-1/2$ on triples with intersection two, and zero elsewhere. We replace
its separate pair roles by signed rectangle mixers, and supply binary
frames specifically for this complex construction.

\paragraph{An explicit partition.}
Let $D_{a,b}(n)$ be the ordered disjointness relation between $a$-subsets
and $b$-subsets of $[n]$, for $0\le a,b\le3$. A plan has counts $(C,L,R)$:
rectangle count, source memberships, and target memberships. An empty
relation has counts zero. If $a=0$ or $b=0$, use one complete rectangle
with counts $(1,\binom na,\binom nb)$. Otherwise compare row stars, column
stars, and every split $n=l+(n-l)$, in that order with increasing $l$.
For a split, each admissible pair of sizes $(x,y)$ gives the product
\[
 D_{x,y}(l)\times D_{a-x,b-y}(n-l).
\]
Use the already chosen unweighted plans for both factors; their three
counts multiply coordinatewise and the counts of distinct cases add.
Choose the alternative minimizing $L+R-C$, retaining the first on ties.
This finite recursion partitions every ordered disjoint pair exactly once.
It is a restricted construction search, not a global optimum claim.

Use $D_{3,3}(25)$ for disjoint triples. For each of the $\binom{25}2$
possible common pairs $J$, use $D_{1,1}(23)$ on the other points and
adjoin $J$ to both singleton families. Intersection-two pairs have a
unique such $J$. The computed counts are
\[
 D_{3,3}(25):(C,L,R)=(21051,172476,253609),\qquad
 D_{1,1}(23):(C,L,R)=(43,104,107).
\]
Thus there are $33951$ rectangles. A rectangle with $a$ sources and $b$
targets uses $a+b-1$ scratch roles, giving
\[
 R_{\rm side}=172476+253609-21051
      +300(104+107-43)=455434.
\]
Exhaustive enumeration checks all $3693800$ ordered neighbor pairs.

\paragraph{Signed mixers and arbitrary scratch.}
Give each rectangle input slots $i_1,\ldots,i_a$ and output slots
$o_1,\ldots,o_b$, sharing only $i_1=o_1$, called the pivot.
The invertible mixer $M$ first adds the other input slots into the pivot,
then adds the pivot into the other output slots. Its output-by-input
submatrix is all ones. The inverse reverses these operations and subtracts.
Let $V$ copy scalar sources to their input slots, and let $J$ inject the
output slots into their targets with coefficient $+1/2$ for disjoint
rectangles and $-1/2$ for intersection-two rectangles. The old central
gather and scatter remain $G,R$. The chronological schedule is
\[
 M,\ -J,\ M^{-1},\ -R,\ V,\ G,\ R,\ M,\ J,\ M^{-1},\ -G,\ -V.
\]
Here a minus sign negates the corresponding additive update.
For initial side scratch $z$ and center scratch $c$, the early operations
subtract $JMz+Rc$ and restore $z$. After adding $Vx,Gx$ to the two scratch
banks, the later operations add $JM(z+Vx)+R(c+Gx)$ and restore both banks.
The net change is $(JMV+RG)x$. By the exact partition and its signs, this
equals $x$: it cancels the central coefficient $(|S\cap T|-1)/2$ for
distinct even-intersection pairs, leaves intersection-one pairs zero,
and leaves the diagonal coefficient one. Thus the map is $y\gets y+x$
for arbitrary scratch. The inverse schedule gives $y\gets y-x$.
The three stages still give $(X,Y)\mapsto(-Y,X)$.

\paragraph{Four binary frame choices.}
Work in $F=\mathbb F_2^{25}$ with its ordinary dot product. For a rectangle,
let $\mathcal X$ be its physical $X$-triple family and $\mathcal Y$ its
physical $Y$-triple family. They are the logical source and target families
in forward stages, and are exchanged in the inverse stage. Choose $U$ by
the following rules in order:
\begin{enumerate}
\item If $|\mathcal X|=1$, use its triple line $U=\langle t_X\rangle$.
\item If $|\mathcal Y|=1$, use $U=t_Y^\perp$.
\item For a disjoint rectangle with both sizes at least two, let $A$ be
the union of its physical $X$ triples and take the coordinate space $U=F^A$.
\item Otherwise the rectangle has intersection two: take
$U=\operatorname{span}\{t_X:X\in\mathcal X\}$.
\end{enumerate}
Every case gives a nondegenerate space with
$\langle t_X\rangle\subset U\subset t_Y^\perp$ for all participating
$X,Y$. Moreover each of the following nonzero spaces is nonalternating:
\[
 U,\quad U^\perp,\quad U\cap t_X^\perp,\quad
 U^\perp\cap t_Y^\perp.
\]
Here the last two are precisely the residuals from a source line to $U$
and from $U$ to a target complement.

For the first two cases, residuals are zero, triple lines, triple
complements, or complements of two orthogonal triple lines. The latter
contain a coordinate unit outside at most six supported coordinates.
For the third case, the source union has at least four points, and the
target union also has at least four points and is disjoint from it.
Thus $U\cap t_X^\perp$ contains a unit in $A\setminus X$, while
$U^\perp\cap t_Y^\perp$ contains a unit outside $A\cup Y$; both
$U$ and $U^\perp$ are coordinate spaces. For the fourth case, distinct
source triples share their fixed pair and have dot product zero and norm
one. They are therefore an orthonormal basis for $U$. Removing one source
line leaves the other basis vectors. A target's singleton coordinate lies
outside the union of the source supports and supplies a unit in $U^\perp$.
For a specified target $Y$, the singleton of a different target supplies
a unit in $U^\perp\cap t_Y^\perp$. Such a target exists because both
family sizes are at least two. The inclusions between nondegenerate spaces
make these residuals nondegenerate as well. This proves the required
orthonormal-basis property in both frame directions.

\paragraph{Gate frames and all residuals.}
At stage $j$ retain the original decomposition
$A=F^{\otimes(j-1)}=B\mathbin\perp P$ and future line $Q$.
Suppressing $Q$, write
\[
 D_0=B\otimes F,\quad D_1=A\otimes F,\quad
 X_t=(B\otimes F)\mathbin\perp(P\otimes\langle t\rangle),
\]
\[
 Y_t=(B\otimes F)\mathbin\perp(P\otimes t^\perp),\qquad
 D_U=(B\otimes F)\mathbin\perp(P\otimes U).
\]
The common frames of the twelve forward operations are respectively
\[
 D_0,D_0,D_0,D_0,X_t,D_1,D_0,D_U,Y_t,D_1,D_1,D_1.
\]
For the inverse stage, reverse and invert the scalar operations, but assign
their physical frames in chronological order as
\[
 D_0,D_0,D_0,X_t,D_U,D_1,D_0,Y_t,D_1,D_1,D_1,D_1.
\]
Each mixer is one gate on the rectangle's entire scratch bank, including
slots it leaves unchanged. Copy and injection gates are grouped by their
data triple. Tensor every displayed frame by $Q$.

Input-only roles follow $D_0\subset X_t\subset D_U\subset D_1$;
output-only roles follow $D_0\subset D_U\subset Y_t\subset D_1$;
the pivot follows both. These statements exchange their physical descriptions
in the inverse stage. The four-space argument above verifies all new
residuals, including direct jumps on roles that skip a copy or injection.
The other factors are the original nonalternating $P,Q,B$ spaces and their
original nonzero residuals. A tensor product of nonalternating spaces contains
a tensor of norm-one vectors; an orthogonal sum with such a summand also
contains one. Hence all nonzero edge residuals admit orthonormal bases.

The physical data boundaries are unchanged. Every central role has exactly
one decrease $D_1\to D_0$ of dimension $h$, so
$L=3v^2h(h+1)=10315500000$. All remaining within-invocation changes are
increasing. Orthogonal weight additivity therefore factors each phase
transition into one signed translation kernel per residual basis vector,
with the same modulo-four signs as the original interface.

\paragraph{Sharing complete auxiliary banks.}
Every auxiliary first gate has label $D_0\otimes Q$ and every last gate
has label $D_1\otimes Q$. Use the disjoint-image permutation $\pi$ from
Section~\ref{sec:complex-stage-reuse}. Match the whole first-stage bank
of invocation $(A,B)$ with the third-stage bank of invocation $(B,\pi(A))$,
keeping each local role index. Their bridge labels are
\[
 E=F\otimes\langle t_A\rangle\otimes\langle t_B\rangle,
 \qquad H=\langle t_B\otimes t_{\pi(A)}\rangle^\perp\otimes F.
\]
They are nested and nondegenerate. A unit
$e_x\otimes e_y\otimes e_q$, with $x\notin B$, $y\notin A$, proves
the bridge residual nonalternating. Each identification saves exactly $m$
residual dimensions, and restoration allows the shared arbitrary inputs.
Consequently
\begin{align*}
 W_{\rm c}^{\rm rect}&=2N+2v^2(R_{\rm side}+h+1)=4843100800000,\\
 s_{\rm c}^{\rm rect}&=W_{\rm c}^{\rm rect}m-2N+2L
                      =75673446297000000,\\
 W_{\rm c}^{\rm rect}m-s_{\rm c}^{\rm rect}&=3703000000,\qquad
 \eta_{\rm c}^{\rm rect}=\frac7{143050000}.
\end{align*}
Every data endpoint and its weight-27 correction is unchanged. Every surviving
scratch role is fixed by the scalar network and still has zero source and full
sink label. The signed bank-exchange and column-tensor corrections therefore
give the original all-role complex phase interface with the new $W,s$.

\begin{proposition}\label{prop:rectangle-complex-interface}
The complex rectangle network supports the original phase interface with
$m=15625$, $W=W_{\rm c}^{\rm rect}$, $s=s_{\rm c}^{\rm rect}$, and
$\sigma=1-5/10^9$.
\end{proposition}
\begin{proof}
The scalar, frame, residual, bridge and endpoint arguments above give the
interface. Exact rational logarithm bounds show
$\eta_{\rm c}^{\rm rect}>(5/10^9)\log m$. The inequality $e^{-x}>1-x$
then gives $s/W<m^{1-5/10^9}$.
\end{proof}

\paragraph{Conservative arithmetic-depth charge.}
The new scalar schedule has at most
$3v^2(4C+4v+4)\le6W$ grouped gates, since $C\le R_{\rm side}$.
Each gate has coefficients in $\{0,\pm1,\pm1/2\}$ and can be implemented
with at most $8W^2$ component additions and binary denominator shifts,
including both real and imaginary components.
The endpoint sign and phase corrections fit within the additional bound
$8(W+m+1)^2$. Thus the old unshared $h=25$ charge
\[
 E_{\rm old}=64(58645352620000+m+1)^3
\]
dominates the total added scalar dependency depth, as also checked exactly.
Translation-kernel child calls are charged separately through $s$.
Since $s<s_{\rm old}<m^5$, the generalized stopped-depth proof may retain
$B_{\rm old}=s_{\rm old}+E_{\rm old}$ and its conservative $C_0$.
The new mixer gates are fixed finite linear maps and introduce no
input-dependent tape count. The compact-control movement and local repair
proofs continue to apply with the new $W,s$.

\section{Latest parameter witness}
Let $a_{\rm b}=296/10^{11}$ and $a_{\rm c}=5/10^9$. Choose
\begin{gather*}
 \tau=1-a_{\rm b},\qquad \sigma=1-a_{\rm c},\qquad
 \epsilon=\frac{1999}{10000},\qquad c=1,\\
 \beta=\frac1{10},\qquad\zeta=\frac1{10000},\qquad
 \delta=\frac1{10^6},\qquad C_1=5-4\beta+\zeta=\frac{46001}{10000},\\
 \lambda=1-\frac{2959}{10^{12}},\qquad
 \lambda'=1-\frac{2958}{10^{12}},\qquad
 \kappa=\frac{59}{10^{11}}.
\end{gather*}
Now $\sigma<\tau$, so the internal exponent is
\[
 \chi=\tau+(1-\beta)\max\{\sigma-\tau,0\}=\tau.
\]
The leaf saving is $(1-\beta)(1-\sigma)=4.5\cdot10^{-9}$,
strictly larger than $1-\lambda'=2.958\cdot10^{-9}$.
Consequently
\[
 \max\{\tau,\sigma,\chi\}<\lambda<\lambda'<1,\qquad
 \max\{\sigma+\beta(1-\sigma),1-c,0\}<\lambda'.
\]
The guard condition is $\epsilon C_1=91955999/10^8<1$.
The reservation estimate is valid for every fixed $c>0$ and uses
$\max\{1-c,0\}$, so $c=1$ is allowed. It reserves only $O(\log p)$
axes; the full spacing is $K=d=\Theta(p^{1999/10000})$.
Since $\ell=\Theta(p^{8001/10000})$, we still have $K=o(\ell)$
and $K/\log p\to\infty$.

All seven assembly margins are
\[
\begin{array}{c|c|c}
 &\text{expression}&\text{value}\\\hline
g_1&1-\epsilon(1+c)&3001/5000\\
g_2&\epsilon c(1-\tau)&73963/(125\cdot10^{12})\\
g_3&\epsilon(1-\lambda')&2956521/(5\cdot10^{15})\\
g_4&(1-\tau)(1-\epsilon)&296037/(125\cdot10^{12})\\
g_5&1/4-\delta-5\epsilon/4&31/250000\\
g_6&1-\delta-\epsilon&800099/10^6\\
g_7&\epsilon&1999/10000
\end{array}
\]
Therefore
\[
 G=\min_i g_i=\frac{2956521}{5\cdot10^{15}}
  =5.913042\cdot10^{-10},\qquad
 G-\kappa=\frac{6521}{5\cdot10^{15}}>0.
\]
The strict positive gap absorbs all retained fixed logarithmic factors.
The unchanged $\epsilon,\delta$ retain the Gaussian and prime comparisons:
$\alpha=\Theta(p^{11999/40000})$,
$\gamma=O(p^{15997/20000})=o(p)$, and $1-2\epsilon>0$.
The guard is $O(d^{46001/10000})=o(p)$, record size remains superpolynomial,
and setup and local repair retain their previous charges. Hence
\[
 T(n)=O\!\left(n(\log n)^{1-59/10^{11}}\right),\qquad
 \frac{59}{10^{11}}>2^{-31},
\]
conditionally on the identified upstream and extension proofs.

\paragraph{What limits this witness.}
The complex saving is now stronger than the retained bit saving.
The constraints $\lambda'>\lambda>\tau$ and the Gaussian margin
$\epsilon<1/5$ imply the scoped ceiling
\[
 \kappa<\epsilon(1-\lambda')<\frac{1-\tau}{5}
    =5.92\cdot10^{-10}<2^{-30}.
\]
The new witness reaches about $99.66\%$ of this ceiling.
It applies only to the fixed bit exponent and these assembly inequalities;
it is not a ceiling for other bit networks or multiplication algorithms.

\section{Verification and reproducibility}
The additional artifacts are \path{scripts/complex_rectangles.py},
\path{certificates/complex-rectangles.json},
\path{tests/test_complex_rectangles.py}, and
\path{patches/complex-rectangles-31.patch}. Run \texttt{make verify}.
The latest patch applies independently to the unmodified pinned manuscript.
The current note is regenerated in place by
\path{scripts/make_complex_rectangles_patch.py}.

Exact checks cover all $3693800$ ordered neighbor pairs, both frame
directions on all $33951$ rectangles, every local source-to-frame and
frame-to-target residual, signed scalar restoration on every input basis
vector of a small rational instance, a complete three-stage shared-scratch
network over an odd field, and actual orthonormal-basis phase factorizations
modulo four. A negative test detects the alternating plane produced by
naively enlarging a triple line to its three-coordinate space. All
parameter comparisons and depth bounds use exact fractions and integers.

The earlier artifacts and certificates are preserved. This is an additional
unreviewed research argument; passing these checks does not establish the
full upstream theorem. The larger exponent saving is not a claim of a
comparable practical runtime improvement.
\end{document}
